\scanpage{24}
\mspage{11}

que cette définition de $C_I$ donne une conclusion aussi forte que $\widetilde u\in C$ pour \uline{tt} couple $I'\subset I''$ de sous-ens. ouverts de $I$.

C’est finalement sous cette forme qu’on va prouver la prop.~2 --- c’est \uline{\unsure{ça|là}} à présent qui me semble \uline{le} bon énoncé dans cette direction. Avec le complément, bien sûr, que si de plus $u_\bullet\in W_{I^\circ}$ i.e. les $u_i$ sont $\in W$, alors $u_\bullet$ est $\unsure{=}{}\in(TC)_I$, \unsure{i.e.|voire} aussi bien que \deleted{\unsure{les}} les \uline{$u_{I'}$ sont $\in TC$} (d’où il \unsure{résultera|résulte} \deleted{\unsure{que}} donc que $\widetilde u$ du cas p.~10 est aussi dans $TC$).

On va procéder par construction transfinie. On peut supposer bien sûr que $I''=I$. On va \inserted{peut $=$} d’abord prouver le cas particulier \deleted{\unsure{où}} \deleted{il existe} où $I'=\varnothing$, i.e. prouver \deleted{que} \unsure{montrer|aussi} que $u_I\in C$. On \deleted{va proc\unsure{.}} construit par réc. transfinie une suite \inserted{str.} croissante d’ouverts

\[
 I_\alpha\subset I,
\]

aboutissant à $I_{\alpha_0}=I$,

en posant

\[
\left\{\begin{aligned}
 I_0&=\varnothing\\
 I_\alpha&=\bigcup_{\beta<\unsure{\alpha|2}}I_{\unsure{\alpha|\beta}}
       &&\text{si }\alpha\text{ ordinal limite}\\
 I_{\alpha\unsure{+|{}}1}&=I_\alpha\cup I_{/s_\alpha}
       &&\text{si }I_\alpha\ne I,\text{ et en choisissant}\\
 &&&s_\alpha\in I\smallsetminus I_\alpha.
\end{aligned}\right.
\]
% archive sheet number: 23

